% GENERATED FILE. Edit canonical lessons, then run npm run generate:books.
% canonical-source-hash: fnv1a64:0bf62c5b0591acbf
% canonical-lessons: HI-C96-tees-chaalees, HI-C96-pachaas-saath, HI-C96-sattar-assi-nabbe, HI-C96-repaso-ginti, HI-C96-sintesis-daam

\chapter{Counting Past Twenty}
\label{ch:hi-ginti}
\hlchaptermodality{104}

\begin{chapteropening}
\textbf{By the end of this chapter:} \emph{I can count the tens from twenty to ninety, and say a price out loud.}

Say pachaas rupaye, then tiis, chaaliis, sattar and sau rupaye, and say which of the two words carries the number.
\end{chapteropening}

\section[tīs cālīs]{\hi{तीस} \hi{चालीस} — two rungs past where the counting stopped}

\label{lesson:HI-C96-tees-chaalees}



\begin{quote}
Count from \textbf{\hi{ग्यारह}} to \textbf{\hi{बीस}}. Then stop, because that is where your counting has ended until now.
\end{quote}

\subsection*{You'll want to know: \hi{तीस}, \hi{चालीस}}
\textbf{\hi{तीस}} (\emph{tīs}) — thirty.

\textbf{\hi{चालीस}} (\emph{cālīs}) — forty.

\begin{grammarlens}[title={Grammar Lens: the tens are the part worth learning}]
\noindent
\begin{tabularx}{\linewidth}{@{}>{\raggedright\arraybackslash}X>{\raggedright\arraybackslash}X@{}}
\toprule
\textbf{number} & \textbf{word} \\
\midrule
20 & \hi{बीस} \\
30 & \hi{तीस} \\
40 & \hi{चालीस} \\
\bottomrule
\end{tabularx}

Say those three in a row and the family resemblance is audible: \textbf{–\hi{ीस}} at the end of all of them. It is the same ending, worn down from the same Sanskrit word for ten, and it will carry you as far as fifty.

\textbf{\hi{तीन}} is three and \textbf{\hi{तीस}} is thirty. \textbf{\hi{चार}} is four and \textbf{\hi{चालीस}} is forty. The openings are not identical to the small numbers, but they are close enough to hold on to, and that is the most help Hindi gives with numbers at this size.
\end{grammarlens}

\subsection*{Your turn}
Say these aloud:

\begin{itemize}
\raggedright
  \item \hi{बीस}, \hi{तीस}, \hi{चालीस} — three times, as a run
  \item the small number each of the two new words leans on
  \item the ending all three share
\end{itemize}

\begin{tcolorbox}[breakable,colback=teal!4,colframe=teal!35!black,title={Before you move on}]
Which number is \textbf{\hi{तीस}}? Which is \textbf{\hi{चालीस}}? Which three of them so far end in the same sound?
\end{tcolorbox}
\section[pacās sāṭh]{\hi{पचास} \hi{साठ} — where the pattern stops helping}

\label{lesson:HI-C96-pachaas-saath}



\begin{quote}
Say \textbf{\hi{बीस}, \hi{तीस}, \hi{चालीस}}. Listen to the ending they share.
\end{quote}

\subsection*{You'll want to know: \hi{पचास}, \hi{साठ}}
\textbf{\hi{पचास}} (\emph{pacās}) — fifty.

\textbf{\hi{साठ}} (\emph{sāṭh}) — sixty.

\begin{grammarlens}[title={Grammar Lens: one keeps the ending, one drops it}]
\noindent
\begin{tabularx}{\linewidth}{@{}>{\raggedright\arraybackslash}X>{\raggedright\arraybackslash}X>{\raggedright\arraybackslash}X@{}}
\toprule
\textbf{number} & \textbf{word} & \textbf{the shared ending} \\
\midrule
40 & \hi{चालीस} & yes \\
50 & \hi{पचास} & a worn version \\
60 & \hi{साठ} & gone \\
\bottomrule
\end{tabularx}

\textbf{\hi{पचास}} still leans on \textbf{\hi{पाँच}}, the five you have counted since the opening chapters, and it still carries a trace of that ending. \textbf{\hi{साठ}} leans on \textbf{\hi{छह}} and keeps nothing else.

That is the shape of the rest of this chapter: the words get less helpful as the numbers get bigger. Knowing that in advance is worth more than a rule that would not hold.

\textbf{\hi{पचास}} is the one to hold on to hardest. It is the price on a tea stall sign and the age of somebody's father, and it is the most common of the four so far.
\end{grammarlens}

\subsection*{Your turn}
Say these aloud:

\begin{itemize}
\raggedright
  \item \hi{बीस}, \hi{तीस}, \hi{चालीस}, \hi{पचास}, \hi{साठ} — the whole run
  \item which of the five keeps least of the shared ending
  \item the small number \hi{पचास} leans on
\end{itemize}

\begin{tcolorbox}[breakable,colback=teal!4,colframe=teal!35!black,title={Before you move on}]
Which number is \textbf{\hi{पचास}}? Which is \textbf{\hi{साठ}}? Which of the two still sounds like the ending you have been hearing?
\end{tcolorbox}
\section[sattar assī nabbe]{\hi{सत्तर} \hi{अस्सी} \hi{नब्बे} — the last three rungs}

\label{lesson:HI-C96-sattar-assi-nabbe}



\begin{quote}
Say \textbf{\hi{पचास}} and \textbf{\hi{साठ}}. One of them kept the old ending and one dropped it.
\end{quote}

\subsection*{You'll want to know: \hi{सत्तर}, \hi{अस्सी}, \hi{नब्बे}}
\textbf{\hi{सत्तर}} (\emph{sattar}) — seventy.

\textbf{\hi{अस्सी}} (\emph{assī}) — eighty.

\textbf{\hi{नब्बे}} (\emph{nabbe}) — ninety.

\begin{grammarlens}[title={Grammar Lens: each one leans on a number you already count}]
\noindent
\begin{tabularx}{\linewidth}{@{}>{\raggedright\arraybackslash}X>{\raggedright\arraybackslash}X@{}}
\toprule
\textbf{small number} & \textbf{ten times it} \\
\midrule
\hi{सात} & \hi{सत्तर} \\
\hi{आठ} & \hi{अस्सी} \\
\hi{नौ} & \hi{नब्बे} \\
\bottomrule
\end{tabularx}

Read that table down the left and you are counting seven, eight, nine — which you have done since early on. Read it down the right and you are counting seventy, eighty, ninety.

\textbf{The opening consonant survives in all three}, and that is the whole of the help on offer. Beyond the first sound, each word has been worn into its own shape and has to be met as a word.

Say the pairs across, one after another: \emph{sāt, sattar. āṭh, assī. nau, nabbe.} Hearing them together is what makes them stick.
\end{grammarlens}

\subsection*{Your turn}
Say these aloud:

\begin{itemize}
\raggedright
  \item \hi{सात} — \hi{सत्तर}, \hi{आठ} — \hi{अस्सी}, \hi{नौ} — \hi{नब्बे}, as three pairs
  \item the whole run from \hi{बीस} to \hi{नब्बे}
  \item what the three new words keep from their small numbers
\end{itemize}

\begin{tcolorbox}[breakable,colback=teal!4,colframe=teal!35!black,title={Before you move on}]
Which number is \textbf{\hi{सत्तर}}? Which is \textbf{\hi{अस्सी}}? Which is \textbf{\hi{नब्बे}}? Which small number does each one lean on?
\end{tcolorbox}
\section[gintī kī tālikā]{\hi{गिनती} \hi{की} \hi{तालिका} — eight rungs, and what sits between them}

\label{lesson:HI-C96-repaso-ginti}



\begin{quote}
Say the last of the three words you met, the one for ninety.
\end{quote}

\begin{grammarlens}[title={Grammar Lens: the ladder}]
\noindent
\begin{tabularx}{\linewidth}{@{}>{\raggedright\arraybackslash}X>{\raggedright\arraybackslash}X@{}}
\toprule
\textbf{number} & \textbf{word} \\
\midrule
20 & \hi{बीस} \\
30 & \hi{तीस} \\
40 & \hi{चालीस} \\
50 & \hi{पचास} \\
60 & \hi{साठ} \\
70 & \hi{सत्तर} \\
80 & \hi{अस्सी} \\
90 & \hi{नब्बे} \\
\bottomrule
\end{tabularx}

Eight words, on top of the hundred you already say. Every price, every age and every house number you are likely to meet sits on or between them.
\end{grammarlens}

\begin{grammarlens}[title={Grammar Lens: what sits between them}]
Here is the honest part, and it is better to hear it than to guess wrong.

\textbf{Hindi puts the small number first.} Twenty-five is \emph{five-twenty}: \textbf{\hi{पच्चीस}} (\emph{paccīs}) opens with the sound of \textbf{\hi{पाँच}}. Twenty-two is \textbf{\hi{बाईस}} (\emph{bāīs}) and opens with \textbf{\hi{दो}}'s vowel.

That order is the opposite of English, and it is the one thing about the in-between numbers a learner can rely on.

What a learner cannot rely on is reassembling them. \textbf{\hi{पच्चीस}} is not \textbf{\hi{पाँच}} and \textbf{\hi{बीस}} stuck together in any way you could work out from the pieces — the two halves have worn into each other. Each of the numbers between the rungs is its own word, and they are met one at a time, over a long while.

So: \textbf{recognise} the shape, know the small number is in front, and learn the rungs cold. That is the part that pays.
\end{grammarlens}

\subsection*{Your turn}
Say these aloud:

\begin{itemize}
\raggedright
  \item the eight rungs in order, without stopping
  \item them backwards, from \hi{नब्बे} down to \hi{बीस}
  \item which end of \hi{पच्चीस} holds the small number
\end{itemize}

\begin{tcolorbox}[breakable,colback=teal!4,colframe=teal!35!black,title={Before you move on}]
How many rungs did this chapter add? Which comes first in \textbf{\hi{पच्चीस}}, the five or the twenty? Can you build twenty-five out of its two pieces? (No — it has to be learned as a word.)
\end{tcolorbox}
\section[Practice]{\hi{पचास} \hi{रुपये} — a price, said out loud}

\label{lesson:HI-C96-sintesis-daam}



\begin{quote}
Say the word for a rupee. You learned where its name came from — a piece of metal that has been given a shape.
\end{quote}

\begin{grammarlens}[title={Grammar Lens: put the two together}]
\textbf{\hi{पचास} \hi{रुपये}} — \emph{fifty rupees}.

Number first, then the thing. That is the whole construction, and you have both halves already: the number arrived in this chapter and the rupee arrived a long way back.

\textbf{\hi{रुपये}} rather than \textbf{\hi{रुपया}}, because more than one is being counted. That bend is the same one the track has been making since the adjectives, and it costs nothing new here.
\end{grammarlens}

\begin{grammarlens}[title={Grammar Lens: the prices you can now read}]
\noindent
\begin{tabularx}{\linewidth}{@{}>{\raggedright\arraybackslash}X>{\raggedright\arraybackslash}X@{}}
\toprule
\textbf{said} & \textbf{meaning} \\
\midrule
\hi{तीस} \hi{रुपये} & thirty rupees \\
\hi{चालीस} \hi{रुपये} & forty rupees \\
\hi{पचास} \hi{रुपये} & fifty rupees \\
\hi{सत्तर} \hi{रुपये} & seventy rupees \\
\hi{सौ} \hi{रुपये} & a hundred rupees \\
\bottomrule
\end{tabularx}

Five prices, and between them most of what a tea stall, a bus fare or a bag of vegetables will ask of you.

\textbf{This is what the counting was for.} A number on its own is a recitation; a number in front of a thing is a price, an age or an address, and that is where the nine rungs stop being a list and start being useful.
\end{grammarlens}

\subsection*{Your turn}
\begin{itemize}
\raggedright
  \item \emph{Say it:} \hi{पचास} \hi{रुपये}
  \item \emph{Say it:} \hi{तीस} \hi{रुपये}, \hi{चालीस} \hi{रुपये}, \hi{सत्तर} \hi{रुपये}, then \hi{सौ} \hi{रुपये}
  \item \emph{Say it:} why the last word is \hi{रुपये} rather than \hi{रुपया}
  \item \emph{Listen:} for the number first and the thing second
\end{itemize}

\begin{tcolorbox}[breakable,colback=teal!4,colframe=teal!35!black,title={Before you move on}]
How do you say \emph{fifty rupees}? Which comes first, the number or the thing? Which form does the rupee take when it is counted?
\end{tcolorbox}
